%%
%% 2020 - 08 -03 φ
%%

\subsubsection{Wahrscheinlichkeitsverteilungen ...}\index{Wahrscheinlichkeitsverteilung}

Für ein Bernoulli-Experiment können wir uns auch Balkendiagramm\index{Balkendiagramm} aufzeichnen.

Nehmen wir wieder unseren Basketballspieler «Basil», der mit Wahrscheinlichkeit von 85\% jeweils trifft. Vorhin hatten wir uns gefragt, wie groß die Wahrscheinlichkeit ist, dass er in 20 Würfen genau 17 Mal [bzw. maximal 17 Mal] trifft.

Dies können wir in einer Wertetabelle und somit auch in einem Balkendiagramm darstellen.
\newpage


\subsubsection*{... genaue Anzahl Treffer}
Notieren wir uns, wie groß die Wahrscheinlichkeit ist, dass Basil \textbf{genau} $n$ Mal trifft innerhalb von 20 Würfen:

\TRAINER{\begin{tabular}{c|cccccccccccccccc}
  n & 0 & 1 & ... & 8 & 9 & 10   & 11   & 12   & 13  & 14  & 15 & 16 & 17 & 18 & 19 &  20\\
  \%& 0 & 0 & ... & 0 & 0 & 0.02 & 0.11 & 0.46 & 1.6 & 4.5 & 10 & 18 & 24 & 23 & 14 & 3.9
\end{tabular}}%% end TRAINER
\noTRAINER{\begin{bbwFillInTabular}{c|cccccccccccccccc}
  n & 0 & 1 & ... & 8 & 9 & 10   & 11   & 12   & 13  & 14  & 15 & 16 & 17 & 18 & 19 &  20\\
  \%& 0 & 0 & ... & 0 &   &      &      &      &     &     &    &    &    &    &    &    
\end{bbwFillInTabular}}%% end no TRAINER


Wenn wir diese Zahlen in ein Balkendiagramm eintragen, so erhalten wir
eine (diskrete) Wahrscheinlichkeitsverteilung:


\bbwCenterGraphic{10cm}{geso/stoch/img/BernoulliSingleBasil.png}

\TNTeop{Wie geht das mit dem TR?
MITSCHREIBEN! (Wir füllen die Tabelle von 10 bis 20 zusammen aus)

DATA -> (4): clear all

DATA -> (rechts) OPS: (3): Sequence (L1)

EXPR IN x:x\\
START x:10\\
END   x:20\\
STEP-SIZE: 1\\
  -> SEQUENCE FILL\\

dann bidomialpdf:

stat-reg/distr (2nd data)\\
-> DISTR (4): binomialpdf\\
wähle LIST\\
TRIALS=n=20\\
p(SUCCESS)=17/20\\
ENTER\\
xLIST: L1\\
SAVE TO: L2\\
CALC (ENTER)

Nun die Werte in L2 ablesen


}%% end TNTeop


\subsubsection*{... kumulierte Anzahl Treffer}
Wie im vorangehenden Kapitel können wir uns die Wahrscheinlichkeiten
auch zusammenzählen und \zB fragen, wie groß ist die
Wahrscheinlichkeit, dass Basil Ballisti (er trifft mit 85\%
Wahrscheinlichkeit) bei 20 Würfen \textbf{maximal} 17
Treffer landet. Dabei darf er für maximal 17 Treffer natürlich auch
weniger Treffer landen. So werden für ein Anzahl Treffer alle
Wahrscheinlichkeiten der weniger erzielten Treffer hinzugezählt. So
erklären sich auch die 100\% um maximal 20 Treffer in 20 Würfen zu
landen.


\TRAINER{
\begin{tabular}{c|cccccccccccccccc}
  n  & 0 & 1 & ... & 8 & 9 & 10  & 11   & 12   & 13  & 14  & 15 & 16 &  17 & 18 & 19 &  20\\
  \% & 0 & 0 & ... & 0 & 0 &  0.02  &0.13   & 0.59 & 2.2 & 6.7 & 17 & 35 &  60 & 82 & 96 & 100
\end{tabular}}%% end TRAINER

\noTRAINER{
\begin{tabular}{c|cccccccccccccccc}
  n  & 0 & 1 & ... & 8 & 9 & 10  & 11   & 12   & 13  & 14  & 15 & 16 &  17 & 18 & 19 &  20\\
  \% & 0 & 0 & ... & 0 & 0 &     &      &      &     &     &    &    &     &    &    &     
\end{tabular}}%% end noTRAINER

Das entsprechende Diagramm sieht wie folgt aus:

\bbwCenterGraphic{10cm}{geso/stoch/img/BernoulliCumulativeBasil.png}
\newpage

Der Vollständigkeit halber hier die Verteilungen unseres
Würfelexperimentes (grün: Normaler Würfel; orange: Gezinkter Würfel):


\bbwCenterGraphic{17cm}{geso/stoch/img/wuerfelverteilung.png}

\newpage
\subsection*{Aufgaben}

\olatLinkGESOKompendium{5.5.4}{50}{29.}

\aufgabenFarbe{Erstellen Sie mit Hilfe des Taschenrechners eine Wertetabelle für die
Wahrscheinlichkeiten, in 10 Würfen 0-mal, 1-mal, 2-mal, ... , 10-mal eine Sechs zu
würfeln.
\\
Erstellen Sie mit diesen Daten ein Balkendiagramm.\vspace{3mm}
\\
Erstellen Sie zusätzlich eine Wertetabelle mit den kumulierten Wahrscheinlichkeiten:
Wie groß ist die Wahrscheinlichkeit mit einem fairen Würfel in 10 Würfen \textbf{maximal} 0-mal, 1-mal, 2-mal, ..., 10-mal eine Sechs zu würfeln. Fürs Verständnis: «maximal 3-mal» heißt «weniger als 4-mal». Erstellen Sie auch hierzu ein Balkendiagramm.
}%% end Aufgabenfarbe

\TNTeop{}

\aufgabenFarbe{Erstellen Sie einen Ausschnitt des «Histogramms» für
unser Pharma-Unternehmen, bei dem wir ein Medikament, das zu 80\,\%
Linderung verspricht an 100 Personen testen.\\
Skizzieren Sie den Ausschnitt von $k$ = 76 bis $k$=84 Personen, indem
Sie jeweils für $k$ angeben, wie groß die Wahrscheinlichkeit ist,
dass \textbf{genau} $k$ Personen Linderung verspüren werden.\\
In $y$-Richtung brauchen Sie Werte von 0 bis 10\,\%.
\\
Für welches $k$ ist diese Wahrscheinlichkeit am größten?
}

\TNTeop{
78: 8.49\,\%

79: 9.46\,\%

80: 9.93\,\%

81: 9.81\,\%

Am größten also bei $k$=80
}%
